\section{Método de Gauss}
\label{sec:gauss}

\subsection*{Combinaciones e independencia de filas}

Una combinación lineal de las filas $F_1,\ldots,F_m$ tiene la forma
$\sum_{i=1}^{m}\lambda_iF_i$, con $\lambda_i\in\K$.
Las filas son \textbf{linealmente independientes} cuando
\[
 \sum_{i=1}^{m}\lambda_iF_i=0
 \quad\Longrightarrow\quad \lambda_1=\cdots=\lambda_m=0.
\]
Son \textbf{dependientes} si existe una relación con algún coeficiente
no nulo; equivale a que alguna fila sea combinación de las demás.
Un conjunto que contiene una fila nula es dependiente.

Para seleccionar filas independientes, se descartan las nulas y se conserva
una fila nueva solo si no es combinación de las ya elegidas. Así se obtiene
un conjunto independiente que permite expresar todas las filas originales.
Si todas son nulas, el conjunto elegido es vacío.

\subsection*{Operaciones elementales y equivalencia}
\label{sec:elementales}

\begin{samepage}
Las operaciones elementales por filas son:
\begin{enumerate}
 \item Intercambiar dos filas distintas: $F_i\leftrightarrow F_j$.
 \item Sumar a una fila un múltiplo de otra:
       $F_i\leftarrow F_i+\lambda F_j$, con $i\ne j$.
 \item Multiplicar una fila por un escalar no nulo:
       $F_i\leftarrow\lambda F_i$, con $\lambda\ne0$.
\end{enumerate}
\end{samepage}
Todas son reversibles: se repite el intercambio, se suma el múltiplo
opuesto o se multiplica por $1/\lambda$, respectivamente.

Dos matrices del mismo orden son \textbf{equivalentes por filas}, $A\ef B$,
si una se obtiene de la otra mediante una sucesión de estas operaciones.
La reversibilidad y la concatenación de sucesiones hacen que esta relación
sea reflexiva, simétrica y transitiva. Las filas de cada matriz se pueden
expresar como combinaciones de las filas de la otra.

Una \textbf{matriz elemental} se obtiene aplicando una operación elemental
a una identidad. Si $A$ tiene $m$ filas y $E$ se obtiene al efectuar
esa operación en $I_m$, el producto $EA$ la efectúa en $A$.
En consecuencia,
\[
 A\ef B\quad\Longleftrightarrow\quad B=E_k\cdots E_1A
\]
para alguna sucesión de matrices elementales; una sucesión vacía deja $A$
igual. El factor situado más a la derecha actúa primero.

\subsection*{Forma escalonada y forma reducida}
\label{sec:escalonada}

El \textbf{pivote} de una fila no nula es su primera entrada no nula,
leída de izquierda a derecha. Una matriz está \textbf{escalonada por filas} si:
\begin{itemize}
 \item Las filas nulas quedan después de todas las no nulas.
 \item El pivote de cada fila no nula está a la derecha del pivote de la anterior.
\end{itemize}
Por ello, debajo de cada pivote solo hay ceros. La forma es
\textbf{escalonada reducida} si, además, cada pivote vale $1$ y es la única
entrada no nula de su columna. Las columnas sin pivote pueden contener
entradas no nulas.

La matriz nula cumple ambas condiciones. Una matriz cuadrada reducida
sin filas nulas es la identidad: debe tener un pivote en cada columna.

\subsection*{Eliminación de Gauss}

El objetivo es obtener una matriz escalonada equivalente por filas.
Se trabaja sobre las filas que aún no tienen pivote:
\begin{enumerate}
 \item Busca la primera columna con una entrada no nula en esas filas.
       Si no existe, la parte restante es nula y el proceso termina.
 \item Intercambia filas para colocar una entrada no nula en la primera
       fila disponible. Esa entrada será el nuevo pivote $a\ne0$.
 \item Para cada entrada $b$ debajo de él, aplica
       $F_i\leftarrow F_i-(b/a)F_r$, donde $F_r$ es la fila del pivote.
 \item Continúa con las filas inferiores y las columnas a la derecha del pivote.
\end{enumerate}
Los ceros ya obtenidos se conservan porque las filas usadas en cada etapa
tienen ceros a la izquierda del nuevo pivote. El proceso termina tras un
número finito de pivotes. La forma escalonada puede variar según las operaciones.

\subsection*{Gauss--Jordan y unicidad}

Para reducir una matriz escalonada, recorre los pivotes desde el último
al primero. Divide cada fila por su pivote y utiliza esa fila para anular
las entradas que quedan por encima. Se obtiene una matriz escalonada reducida.

\begin{ejemplo}
La cadena siguiente muestra la eliminación y la reducción:
\[
 \begin{pmatrix}1&2&1\\2&4&3\\0&0&2\end{pmatrix}
 \longrightarrow
 \begin{pmatrix}1&2&1\\0&0&1\\0&0&0\end{pmatrix}
 \longrightarrow
 \begin{pmatrix}1&2&0\\0&0&1\\0&0&0\end{pmatrix}.
\]
En la primera flecha se aplican, en este orden,
$F_2\leftarrow F_2-2F_1$ y $F_3\leftarrow F_3-2F_2$;
en la segunda, $F_1\leftarrow F_1-F_2$. Los pivotes están en las columnas
primera y tercera; la segunda no necesita ser nula.
\end{ejemplo}

Toda matriz $A$ tiene una \textbf{única forma escalonada reducida por filas},
denotada $H_f(A)$, y
\[
 A\ef B\quad\Longleftrightarrow\quad H_f(A)=H_f(B).
\]
La unicidad se basa en que dos matrices reducidas equivalentes deben
tener las mismas posiciones de pivote y las mismas filas. No depende
de la elección de operaciones durante el cálculo.

\textbf{Control del cálculo.} Puede elegirse un pivote sencillo para reducir
las fracciones. Cada operación utiliza las filas que existen en ese momento;
si se escriben varias juntas, su orden debe ser inequívoco. Si una entrada
depende de un parámetro, hay que separar el caso en que vale cero antes
de dividir por ella o utilizarla como pivote.

\subsection*{Operaciones por columnas y equivalencia general}

Las operaciones elementales por columnas se obtienen sustituyendo
$F_i$ por $C_i$ en las operaciones por filas. Una operación en
$I_n$ produce una matriz elemental $E$ que actúa sobre las columnas
de una matriz $A$ de orden $m\times n$ mediante $AE$.

La equivalencia por columnas y la reducción correspondiente se obtienen
trasponiendo:
\[
 A\ec B\quad\Longleftrightarrow\quad A^t\ef B^t,
 \qquad H_c(A)=\bigl(H_f(A^t)\bigr)^t.
\]
Así, la forma escalonada reducida por columnas también existe y es única.

La \textbf{equivalencia} $A\sim B$ permite combinar operaciones elementales
por filas y por columnas. Equivale a poder escribir
\[
 B=PAQ,
\]
con $P$ y $Q$ productos de matrices elementales de órdenes $m$ y $n$,
respectivamente. También es una relación de equivalencia. La equivalencia
por filas o por columnas implica esta equivalencia general.
